\documentclass[10pt,a4paper]{amsart}
\usepackage[margin=19mm]{geometry}
\usepackage{amsmath,amssymb,mathtools}
\usepackage[scr=rsfs,cal=euler]{mathalfa}
\usepackage{xfrac}
\usepackage[unicode,hidelinks,bookmarks=false]{hyperref}
\newcommand{\ie}{i.e.\ }
\newcommand{\cf}{cf.\ }
\newcommand{\resp}{respectively\ }
\newcommand{\Subsection}{\subsection}
\providecommand{\nobreakdash}{\nobreak\hbox{-}}
\setlength{\emergencystretch}{2em}
\multlinegap0pt
\usepackage{amssymb}
\usepackage{float}
\usepackage{tikz}
\newtheorem{theorem}{Theorem}
\newcommand{\N}{\mathbb{N}}
\newcommand{\R}{\mathbb{R}}
\title{Hitting sets and colorings of hypergraphs}
\author{Bal\'azs Bursics, Bence Csonka, Luca Szepessy}
\date{Source: arXiv:2307.12154v3 (CC BY 4.0)}
\begin{document}
\maketitle

\noindent\emph{Source and licence.} Hitting sets and colorings of hypergraphs. Version arXiv:2307.12154v3 (CC BY 4.0), distributed by arXiv under the Creative Commons Attribution 4.0 International licence, http://creativecommons.org/licenses/by/4.0/, as recorded in the arXiv licence metadata for this exact version and shown on the abstract page https://arxiv.org/abs/2307.12154v3. Full licence notice: \texttt{licenses/arxiv-2307.12154-CC-BY-4.0.txt}.\par\smallskip

\noindent\textit{Editorial scope.} Counted unit: the article's theorem vegtelen\_sorozat (source0.tex lines 320--342). The article's complete original proof of it is delivered with it (source0.tex lines 687--712, one \texttt{proof} environment of 26 lines, which the article places away from the statement). No mathematics is written, repaired or re-derived: the statement and the proof are the article's own lines, in the article's own notation, and no retained line is substituted or re-flowed.  No further source-local context is needed: every cross-reference of every retained line resolves inside the two delivered spans.  One declared adaptation: exactly 1 ancillary strings inside the retained spans are omitted (image-inclusion commands and, where the source repeats a label, the repeated label definitions) and no other line differs from the source; the figure environments, with their placement, captions and labels, are kept, and the package records the omitted strings in fidelity receipts bound to the affected bodies.  No retained line contains a citation, so the wrapper carries no bibliography and no bibliography entry.  The retained lines contain the article's own diagram code, which the wrapper typesets with the standard tikz packages; no image file is delivered and the package declares no asset dependency.  Editorial formatting only, in addition: an AMS \texttt{amsart} preamble that restates the article's own declarations and macros used by the retained lines, namely the environments \texttt{theorem} on separate counters, together with the standard packages those lines need.  The retained environments print in this document as Theorem 1 in order of appearance, whereas the article prints the same items with the numbering of its own full document.  The complete native source of the article as distributed in the arXiv e-print for this exact version is bundled unchanged as \texttt{sources/144655/source0.tex}.
\begin{theorem}\label{vegtelen_sorozat}

\begin{enumerate}
    \item 
    Let $D := \{t^i:i \in \mathbb{N}\}$ for some $t \in \mathbb{N} \setminus \{0,1\}$, then
    \[
    m_{D, \mathbb{N}}^{\infty}(k) \leq m_{\mathcal{T}_3}(k).
    \]
    More generally, this also holds if $D = \{d_i: i\in\N, d_i \in \mathbb{N}^+, d_{i-1} | d_{i}\}$, $d_0:=1$.
    \item 
    Let $p,q$ be positive coprime integers, and suppose $M$ contains at most one element of every residue class modulo $pq$, then
    \[
    m_{\{p^iq^j:i,j \in \mathbb{N}\}, M}^{\infty}(k) \leq \max(p,q)\cdot m_{\mathcal{T}_3}.
    \]
    Moreover, if $|M|=1$, then $m_{\{p^iq^j:i,j \in \mathbb{N}\}, M}^{\infty}(k) = m_{\mathcal{T}_3}$.
    \item 
    Let $p_1, p_2, p_3$ be positive coprime integers, then 
    \[
    m_{\left\{ p_1^{j_1}p_2^{j_2}p_3^{j_3}:j_1,j_2,j_3 \in \mathbb{N}\right\},\{0\}}^{\infty}(k) = \infty.
    \]

\end{enumerate}
\end{theorem}

\begin{proof}[Proof of Theorem~\ref{vegtelen_sorozat}. (1):]
    We will show that for every $A \in \mathcal{A}_{D,\N}^{\infty}$ there exists $T_A \in \mathcal{T}_3$ and $\phi :V(A)\to V(T_A)$ injection such that for every hyperedge $E \subseteq V(A)$ the set $\phi [E]\subseteq V(T_A)$ is also a hyperedge of $T_A$. Therefore, if we can $k$-color $\left(T_A\right)_m$ properly for $m:=m_{\mathcal{T}_3}(k)$, we obtain a polychromatic $k$-coloring of $A_{m}$.
    
    The vertex set of $A \in \mathcal{A}_{D,\N}^{\infty}$ contains finitely many natural numbers. We will assign points in $\mathbb{R}^3$ to them as the vertices of $T_A$ such that for every hyperedge $E \subseteq V(A)$ the set $\phi [E]\subseteq V(T_A)$ can be defined as the intersection of an octant and $V(T_A)$. Fix the difference set D. For the sake of simplicity, we will assign a point in $\mathbb{R}^3$ to every natural number, such that every subset of $\mathbb{N}$ which could be a hyperedge in a hypergraph $A \in \mathcal{A}_{D,\N}^{\infty}$ can be defined by an octant in $\mathbb{R}^3$. 
    
    Firstly, let us discuss the case when $D = \{t^i: i \in \mathbb{N}\}$ for some $t \in \mathbb{N}^+$. We will place axis-parallel squares recursively into each other on the $y$-$z$ plane, each of them corresponding to a natural number. This natural number is placed in the bottom left corner of the square, see Figure \ref{negyzetek}.
    
    \begin{figure} [h!]
    \centering
    
    \caption{Placement of the squares perpendicular to the $x$-axis}
    \label{negyzetek}
\end{figure}
    
    The $0$th generation of the recursion is a square $S_0$ in the $y$-$z$ plane, corresponding to $0$, say the unit square $[0,1]\times[0,1]$. Place a point $P_0$ in the southwest corner of it. For the first generation, take those natural numbers whose $t$-base form contains only one non-zero bit. Accordingly, the first generation consists of squares $S_1, S_2,...,S_{t-1}, S_{10}, S_{20}, ..., S_{[t-1]0}, S_{100}, S_{200}, ...$ (with the indices being in base $t$) corresponding to $1, 2, ..., t-1, t, 2t, ..., (t-1)t, t^2, 2t^2, ...$, respectively. We place them diagonally (from northwest to southeast) into $S_0$, according to Figure \ref{negyzetek}, and we place a point $P_a$ into the southwest corner of every $S_a$. In terms of coordinates, the $j^\text{th}$ square of the first generation is $\left[1-\frac{1}{j+1}, 1-\frac{1}{j+2}\right]\times\left[\frac{1}{2}+\frac{1}{j+2},\frac{1}{2}+\frac{1}{j+1}\right]$.
    
    The second generation of squares correspond to natural numbers whose $t$-base form contains two nonzero bits. The value and place of the smallest nonzero bit defines that first-generation square in which we put our second-generation square: in $S_a$ where the value of $a$ is $k_1 \cdot t^{i_1}$, we place the squares corresponding to the values $k_1 \cdot t^{i_1} + k_2 \cdot t^{i_2}$ (where $i_1<i_2$), in a diagonal, ordered by the absolute value of $k_2\cdot t^{i_2}$. Take the similarity transformation that maps $S_0$ onto $S_a$ and preserves all directions. The squares of the second generation are placed as the images of the first-generation squares, in the prescribed order. And so on, the $k$-th generation consists of squares corresponding to numbers whose $t$-base form contains exactly $k$ nonzero bits. Note that a square $S_b$ contains another square $S_c$ exactly when the $t$-base form of $c$ is a finat segment of the $t$-base form of $b$. 
    
    As in the 0th and first generation, into every square $S_a$ placed in the $k$-th generation we put a point $P_a$ to the southwest corner. Now we have points in the $y$-$z$ plane assigned to every $n\in\mathbb{N}$, if $n$ has $k$ pieces of nonzero bits in its $t$-base form, we have placed a square for it in the $k$-th generation. Finally, let $\phi (n)=(n,P_n)\in\R^3.$
    
    We have to prove that for every set of form $E = \{a_0+ i \cdot t^j: i \in \mathbb{N}\}$, $a_0, j \in \mathbb{N}$ there exists an octant which contains exactly these points. Firstly, fix $a_0$ and $j$ so that $a_0<t^j$. If we project back our points to the $y$-$z$ plane, we can almost "cut out" the sequence with an axis-parallel quarter-plane; see Figure \ref{negyzetek}. Take the square $S_b$ corresponding to $a_0$. According to the construction, the square $S_c$ corresponding to $a_0 + t^j$ is contained in $S_b$. Let $e$ be the east bordering line of $S_b$ and $f$ be the north bordering line of $S_c$. The lines $e$ and $f$ define four quarter-planes, and let $Q$ be the southwest one with respect to the orientation given by $y$ and $z$. $Q$ contains points corresponding to numbers that are at least $a_0$ and are divisible by $t^j$, or that are smaller than $a_0$. Notice that every number in form of $n=a_0+k \cdot t^j$, $k\in \mathbb{N}$ is contained in $Q$ since only $S_b$ contains numbers congruent to $a_0$ modulo $t^j$, and points above $f$ in $S_b$ correspond to numbers which are the sum of $a_0$ and $r_j$ where $r_j$ has a $t$-base form with $0$ in its $j^{th}$ bit. 
    
    For the case $a_0 \geq t^j$, we can define some $a_{-1}, a_{-2},...,a_{-k}$ such that $0\leq a_{-k} < t^j$, and $\{a_i\}_{i \in \{-k, ..., -1\} \cup \mathbb{N}}$ is still an arithmetic progression, and we can build the same construction for this extended arithmetic progression as in the case $a_0<t^j$. Take the axis-parallel lines $e$ and $f$ as in the previous case, and assume that they are defined by the equations $y=y_0$ and $z=z_0$, respectively. Then the octant in $\mathbb{R}^3$ corresponding to $E$ is the set $\{(x,y,z): x \geq a_0, y \leq y_0, z \leq z_0\}$. 
    
    The general case when $D = \{d_i: i\in\N, d_i \in \mathbb{N}^+, d_{i-1} | d_{i}\}$, $d_0:=1$ is very similar. We use almost the same construction, we just replace the $t$-base form of natural numbers with the following: Assume that for every $k<n$ we have defined a form $k = \sum\limits_{j=0}^{i-1} c_j \cdot d_j$. For $n \in \mathbb{N}$, let $i$ be the maximal index for which $d_i \leq n$. Let $c_i$ be the maximal natural number for which $c_i \cdot d_i \leq n$. Notice that $c_i<\frac{d_{i+1}}{d_i}$. If $k := n-c_i \cdot d_i = \sum\limits_{j=0}^{i-1} c_j \cdot d_j$, write $n$ in form $n = \sum\limits_{j=0}^{i} c_j \cdot d_j$. We build up the construction in the same way, we just replace the $t$-base form of numbers with the sequence $...c_2c_1c_0$.
    \end{proof}

\end{document}
